% ====================================================================
% ch2v22_notation.tex — Ch2(LCO) 기호표: 계승 2단 (R1 초판 — R3 에서 확정)
%   규약(D22): 계승 기호는 재정의하지 않는다 — 정의는 Chapter 1 해당 위치가 유일 기준.
% ====================================================================
\subsection*{기호 --- 계승과 신규}
본 장의 기호는 두 단으로 관리한다. \textbf{계승}(Chapter 1 정의 그대로 --- 재정의 없음):
전이 지표 $j$·용량 $Q_j$·평형 중심 $U_j(T)$~\eqref{eq:Uj}·폭 $w_j=n_jRT/F$·평형 진행률
$\xi_{\eq,j}$~\eqref{eq:xieq}·상호작용 $\Omega_j$ 와 히스 gap~\eqref{eq:dUhys}·분기 인자
$\gamma_j$~\eqref{eq:Ubranch}·방향 부호 $\sigma_d$·반응 엔트로피 $\Delta S_{\rxn,j}$·추출 분율
$\bar x$ 와 전하 보존 반전~\eqref{eq:sm-mc-balance}·합산식~\eqref{eq:sum}. \textbf{신규}(본 장
도입 --- 각 도입 절이 정의 기준): 탈리튬화$\cdot$충방전 재배선 $\sigma_d$ 슬롯 규약(\S\ref{sec:lco-direction}),
질서상$\cdot$MIT 2상역 파라미터(\S\ref{sec:lco-hys-od}$\cdot$\S\ref{sec:lco-hys-mit}), 전자 엔트로피 $\Delta S_e$$\cdot$상태밀도
게이트 $g(E_F,x)$$\cdot$게이트 중심 $x_\mathrm{MIT}$(\S\ref{sec:lco-electronic}), 슬롯 분해
$\Delta S^\mathrm{config/vib/e}_j$(\S\ref{sec:lco-decomp}).
